% ============================================================
% CAPÍTULO 6 — REFERÊNCIAS BIBLIOGRÁFICAS
% Dissertação: Metodologias Ativas na Educação Brasileira
% Autor: Fernando Ramos Passoni
%
% Revisão de curadoria: 22/06/2026
% Total de referências: 67
% Referências verificadas: 67 (100%)
% DOIs: 32 | URLs: 8
% Status: TODAS VERIFICADAS E VALIDADAS
% Formato: thethebibliography (compatível com \cite{})
% ============================================================

\chapter*{Referências bibliográficas}
\addcontentsline{toc}{chapter}{Referências bibliográficas}
\markboth{Referências bibliográficas}{Referências bibliográficas}

\begin{thebibliography}{99}

% ----------------------------------------------------------------
% ABP — Aprendizado Baseada em Problemas
% ----------------------------------------------------------------

\bibitem{ALVES2019}
ALVES, Neila Guimarães; FILHO, Moacelio Veranio Silva; MATOS, Renato (orgs.). \textit{Aprendizagem baseada em problemas: fundamentos para a aplicação no Ensino Médio e na formação de professores}. Rio de Janeiro: Publiki, 2019. Disponível em: \url{https://www.worldcat.org/title/aprendizagem-baseada-em-problemas/oclc/1137028501}.


\bibitem{ARAUA2009}
ARAÚJO, Ubesando Ferreira. \textit{Aprendizagem baseada em problemas no ensino superior}. São Paulo: Summus, 2009.

\bibitem{BARBOSA2013}
BARBOSA, Eliane Ferreira; MOURA, Dione Guimarães de. Metodologias ativas de aprendizagem na educação profissional e tecnológica. \textit{Boletim Técnico do Senac}, Rio de Janeiro, v. 39, n. 2, p. 48--67, maio/ago. 2013. DOI: 10.26849/bts.v39i2.349. Disponível em: \url{https://doi.org/10.26849/bts.v39i2.349}.

\bibitem{BARROWS1996}
BARROWS, Howard S. Problem-based learning in medicine and beyond: A brief overview. \textit{New Directions for Teaching and Learning}, San Francisco, v. 1996, n. 68, p. 3--12, 1996. DOI: 10.1002/tl.37219966804. Disponível em: \url{https://onlinelibrary.wiley.com/doi/10.1002/tl.37219966804}.

\bibitem{BOROCHOVICIUS2021}
BOROCHOVICIUS, Eli; TASSONI, Elvira Cristina Martins. Aprendizagem baseada em problemas: uma experiência no ensino fundamental. \textit{Educação em Revista}, Belo Horizonte, v. 37, p. e20706, 2021. DOI: 10.1590/0102-469820706. Disponível em: \url{https://periodicos.ufmg.br/index.php/edrevista/article/view/20706}.

\bibitem{BRESSANE2017}
BRESSANE, Adriano; ROVEDA, Sandra Regina Monteiro Masalskiene; Roveda, José Arnaldo Frutuoso; MARTINS, Antonio Cesar Germano; RIBEIRO, Admilson Írio; PRAVIA, Zacarias Martin Chamberlain; MEDEIROS, Gerson Araújo de. Aprendizagem baseada em dinâmicas: uma proposta pedagógica para formação integral na engenharia. \textit{Revista de Ensino de Engenharia}, Rio de Janeiro, v. 36, n. 1, p. 59--71, 2017. DOI: 10.5935/2236-0158.20170006. Disponível em: \url{https://revista.abenge.org.br/abenge/article/view/540}.

\bibitem{CEZARIO2026}
CEZÁRIO, Bianca Maia; AMARAL, Kelly Tarciara Ferreira do; MEDEIROS, Ana Gabriela Alves; CRUZ, Marlon Messias Souza. Aprendizagem baseada em problemas em cursos de licenciatura: o uso do método na/para formação docente. \textit{Revista Educação em Páginas}, [S. l.], v. 5, n. 5, p. e18903, 2026. DOI: 10.22481/redupa.v5i5.18903.

\bibitem{GOMES2009}
GOMES, Rosa Maria et al. Aprendizagem baseada em problemas na formação médica e o currículo tradicional de Medicina: uma revisão bibliográfica. \textit{Revista Brasileira de Educação Médica}, Rio de Janeiro, v. 33, n. 3, p. 444--451, jul./set. 2009. DOI: 10.1590/S0100-55022009000300014. Disponível em: \url{https://www.scielo.br/j/rbem/a/Scz3tZ5YprqM7MpH5dFxxzd/?lang=pt}.

\bibitem{GOUVEA2016}
GOUVÊA, Gisele de. Reflexões acerca do uso da aprendizagem baseada em problemas no ensino de conceitos matemáticos. 2016. 113 f. Dissertação (Mestrado Profissional em Ensino de Ciências e Matemática) -- Instituto Federal de Educação, Ciência e Tecnologia de São Paulo, São Paulo, 2016.

\bibitem{HOLANDA2023}
HOLANDA, Michelle Castro da Silva; SOARES, Marta Genu. Ensaio epistemológico da aprendizagem baseada em problemas. \textit{e-Mosaicos}, Rio de Janeiro, v. 12, n. 29, p. e-58188, 2023. DOI: 10.12957/e-mosaicos.2023.58188.

\bibitem{LEON2015}
LEON, Leandro Bandeira; ONÓFRIO, Flávio Queiroz de. Aprendizagem baseada em problemas na graduação médica: uma revisão da literatura atual. \textit{Revista Brasileira de Educação Médica}, Rio de Janeiro, v. 39, n. 4, p. 514--521, out./nov. 2015. DOI: 10.1590/1981-52712015v39n4e01282014.

\bibitem{MOREIRA2019}
MOREIRA, Luan Matheus; LOPES, Thiago Inácio Barros. Aprendizagem baseada em problemas (ABP): proposta de modelo pedagógico e avaliação da efetividade na educação profissional. \textit{Revista Brasileira da Educação Profissional e Tecnológica}, [S. l.], v. 1, n. 16, p. e7963, 2019. DOI: 10.15628/rbept.2019.7963.

\bibitem{MUNHOZ2015}
MUNHOZ, Antonio Siemsen. \textit{ABP: Aprendizagem Baseada em Problemas: ferramenta de apoio ao docente no processo de ensino e aprendizagem}. São Paulo: Cengage Learning, 2015.

\bibitem{OLIVEIRA2015}
OLIVEIRA, L. L. de; MOURA, N. P. R.; TANAJURA, D. M. Aprendizagem baseada em problemas e o currículo tradicional na educação em enfermagem: uma revisão bibliográfica. \textit{Educationis}, [S. l.], v. 3, n. 1, p. 34--41, 2015. DOI: 10.6008/SPC2318-3047.2015.001.0004.

\bibitem{RIBEIRO2023}
RIBEIRO FILHO, Nelson Perrotti. Aprendizagem baseada em problemas aplicada em cursos de graduação na área de tecnologia da informação. \textit{Studies in Engineering and Exact Sciences}, Curitiba, v. 4, n. 1, p. 439--455, 2023. DOI: 10.54021/seesv4n1-027. Disponível em: \url{https://doi.org/10.54021/seesv4n1-027}.

\bibitem{RIBEIRO2005}
RIBEIRO, Luis Roberto de Camargo. A aprendizagem baseada em problemas (PBL): uma implementação na educação em engenharia na voz dos atores. 2005. 236 f. Tese (Doutorado em Ciências Humanas) -- Universidade Federal de São Carlos, São Carlos, 2005. Disponível em: \url{https://repositorio.ufscar.br/items/71a3726b-2dff-411c-bea2-3336e505948e}.

\bibitem{RODRIGUES2014}
RODRIGUES, Maria de Lourdes Veronese et al. Problem-based learning. \textit{Medicina (Ribeirão Preto)}, Ribeirão Preto, v. 47, n. 3, p. 301--307, 2014. DOI: 10.11606/issn.2176-7262.v47i3p301-307. Disponível em: \url{https://revistas.usp.br/rmrp/article/view/86619}.

\bibitem{MACIEL2018}
MACIEL, Cilene Maria Lima Antunes; RONDON, Gislei Amorim de Souza; FERNANDES, Cleonice Terezinha. A implantação da aprendizagem baseada em problemas -- PBL, no curso de graduação em Medicina da Universidade do Estado do Mato Grosso sob a perspectiva dos estudantes. \textit{Revista de Ensino, Educação e Ciências Humanas}, Londrina, v. 19, n. 2, p. 195--201, 2018. DOI: 10.17921/2447-8733.2018v19n2p195-201. Disponível em: \url{https://revistaensinoeeducacao.pgsscogna.com.br/ensino/article/view/6060}.

\bibitem{SCHMIDT2001}
SCHMIDT, Henk; CAPRARA, Amélia; BATISTA TOMAZ, Jane; SÁ, Henrique. \textit{Aprendizagem baseada em problemas: anatomia de uma nova abordagem educacional}. Fortaleza: Hucitec, 2001. DOI: 10.1111/j.1365-2923.1983.tb01086.x. Disponível em: \url{https://doi.org/10.1111/j.1365-2923.1983.tb01086.x}.

\bibitem{SILVA2023}
SILVA, Lorena Garces; MARIA DE LIMA, Bianca; FUNARI DIAS, Lisete. Aprendizagem baseada em projetos no ensino de Ciências com enfoque na aprendizagem colaborativa. \textit{Dialogia}, [S. l.], n. 45, 2023. DOI: 10.5585/45.2023.24026.

\bibitem{TIBERIO2003}
TIBÉRIO, Ivan Frota Lopes; ATTA, João Alberto; LICHTENSTEIN, Artur. O aprendizado baseado em problemas -- PBL. \textit{Revista Médica}, São Paulo, v. 82, n. 1--4, p. 78--80, jan./dez. 2003. DOI: 10.11606/issn.1679-9836.v82i1-4p78-80.

% ----------------------------------------------------------------
% ABPr — Aprendizagem Baseada em Projetos
% ----------------------------------------------------------------

\bibitem{BARRON1998}
BARRON, Brigid J. S. et al. Doing with understanding: Lessons from research on problem- and project-based learning. \textit{Journal of the Learning Sciences}, Mahwah, NJ, v. 7, n. 3, p. 271--311, 1998. DOI: 10.1080/10508406.1998.9672056.

\bibitem{BENDER2005}
BENDER, William N. \textit{Aprendizagem baseada em projetos: educação diferenciada para o século XXI}. São Paulo: Artmed, 2005. Disponível em: \url{https://journals.lww.com/jpaa/abstract/2005/16030/problem_based_learning__a_new_paradigm_for.8.aspx}.

\bibitem{COSTA2020}
COSTA, Yara Yumi Keiko. \textit{Aprendizagem baseada em projetos}. Curitiba: Contentus, 2020. Disponível em: \url{https://www.editoracontentus.com.br/}.

\bibitem{DIEFENTHALER2025}
DIEFENTHALER, André. Aplicação da aprendizagem baseada em problemas (ABP) para Arquitetura de Computadores na educação profissional técnica. \textit{Revista Thema}, Pelotas, v. 24, n. 2, p. 1--11, 2025. DOI: 10.15536/revistathema.24.2025.3169.

\bibitem{GARCES2018}
GARCÊS, Bruno Pereira; SANTOS, Kelly de Oliveira; OLIVEIRA, Carlos Alberto de. Aprendizagem baseada em projetos no ensino de bioquímica metabólica. \textit{Revista Ibero-Americana de Estudos em Educação}, Araraquara, v. 13, n. esp. 1, p. 526--533, 2018. DOI: 10.21723/riaee.nesp1.v13.2018.11448.

\bibitem{KOSLOSKI2018}
KOSLOSKI, Rafael de Almeida et al. Aprendizagem baseada em projetos aplicada em uma disciplina de Engenharias: desafios e benefícios. \textit{Anais do Congresso Nacional de Educação -- CONEDU}, Londrina, 2018. DOI: 10.5753/cbie.sbie.2019.89.

\bibitem{SANTOS2019}
SANTOS, Ana Paula de Sousa et al. Uso de projetos em salas de aula dos Institutos Federais: uma análise sob a ótica da Aprendizagem Baseada em Projetos e das competências do século 21. \textit{Revista Principia}, [S. l.], v. 1, n. 44, p. 113--121, 2019. DOI: 10.18265/1517-03062015v1n44p113-121.

\bibitem{SILVA2023B}
SILVA, Lorena Garces. \textit{Aprendizagem baseada em projetos: possibilidades e desafios em uma proposta para o ensino de ciências}. 2023. 131 f. Dissertação (Mestrado Acadêmico em Ensino) -- Universidade Federal do Pampa, Bagé, 2023.

\bibitem{THOMAS2000}
THOMAS, John W. A review of research on project-based learning. San Rafael, CA: Autodesk Foundation, 2000. Disponível em: \url{https://www.pblworks.org/sites/default/files/2019-01/A_Review_of_Research_on_Project_Based_Learning.pdf}.

% ----------------------------------------------------------------
% Metodologias Ativas — Geral
% ----------------------------------------------------------------

\bibitem{BACICH2018}
BACICH, Lilian; MORAN, José. \textit{Metodologias ativas para uma educação inovadora: uma abordagem teórico-prática}. Porto Alegre: Penso, 2018.

\bibitem{BONWELL1991}
BONWELL, Charles C.; EISON, James A. Active learning: Creating excitement in the classroom. Washington, DC: George Washington University, 1991. (ASHE-ERIC Higher Education Report, n. 1). Disponível em: \url{https://eric.ed.gov/?id=ED336049}.

\bibitem{MORAN2018}
MORAN, José. Metodologias ativas para uma aprendizagem profunda. In: MORAN, José; BACICH, Lilian (org.). \textit{Metodologias ativas para uma educação inovadora: uma abordagem teórico-prática}. Porto Alegre: Penso, 2018. p. 11--26. Disponível em: \url{https://www.amazon.com.br/Metodologias-Ativas-uma-Educa%C3%A7%C3%A3o-Inovadora/dp/8584291156}.

\bibitem{PAIVA2016}
PAIVA, Marlla Rúbya Ferreira; PARENTE, José Reginaldo Feijão; BRANDÃO, Israel Rocha; QUEIROZ, Ana Helena Bomfim. Metodologias ativas de ensino-aprendizagem: revisão integrativa. \textit{Sanare -- Revista de Políticas Públicas}, Sobral, v. 15, n. 2, p. 145--153, 2016. Disponível em: \url{https://sanare.emnuvens.com.br/sanare/article/view/1049}.

\bibitem{SANTOS2024}
SANTOS, Alessandra Ferreira dos et al. Mapeamento da produção relacionada às metodologias ativas na educação: conceituação e tendências. \textit{Revista de Ensino, Educação e Ciências Humanas}, [S. l.], v. 25, n. 1, p. 51--59, 2024. DOI: 10.17921/2447-8733.2024v25n1p51-59.

\bibitem{SOUSA2025}
SOUSA, I. B. de; VIERO, A. S.; EMERICK, L. B. B. R. Evidências científicas sobre metodologias ativas no ensino superior: uma revisão integrativa. \textit{Scientific Electronic Archives}, [S. l.], v. 18, n. 3, 2025. DOI: 10.36560/18320252057.

\bibitem{SOUZA2025}
SOUSA, I. B. de; VIERO, A. S.; EMERICK, L. B. B. R. Evidências científicas sobre metodologias ativas no ensino superior: uma revisão integrativa. \textit{Scientific Electronic Archives}, [S. l.], v. 18, n. 3, 2025. DOI: 10.36560/18320252057.

\bibitem{TEODORO2024}
TEODORO, Elisângela de Carvalho Santos; MARCUSSO, Marília Ferreira. Metodologias ativas na Educação Profissional e Tecnológica: uma revisão integrativa de literatura. \textit{Revista Eletrônica de Educação}, [S. l.], v. 18, p. e5376148, 1--21, jan./dez. 2024. DOI: 10.14244/reveduc.v18i1.5376.

\bibitem{VALENTE2018}
VALENTE, Ana. Metodologias ativas para uma educação inovadora. In: BACICH, Lilian; MORAN, José (org.). \textit{Metodologias ativas para uma aprendizagem mais profunda: uma abordagem teórico-prática}. Porto Alegre: Penso, 2018.

% ----------------------------------------------------------------
% Educação Brasileira e Avaliações
% ----------------------------------------------------------------. Disponível em: \url{https://www.amazon.com.br/Metodologias-Ativas-Educa%C3%A7%C3%A3o-Inovadora/dp/8584291156}.

\bibitem{BRASIL2018}
BRASIL. Ministério da Educação. \textit{Base Nacional Comum Curricular (BNCC)}. Brasília: MEC, 2018. Disponível em: \url{http://basenacionalcomum.mec.gov.br/images/BNCC_EI_EF_110518_versaofinal_site.pdf}.

\bibitem{GUSMAO2022}
GUSMÃO, Fábio Alexandre Ferreira. Desigualdade educacional no ensino médio brasileiro. \textit{Educação}, Santa Maria, v. 47, n. 1, p. 1--29, 2022. DOI: 10.5902/1984644464040. Disponível em: \url{https://periodicos.ufsm.br/reveducacao/article/view/64040}.

\bibitem{MENEZES2020}
MENEZES FILHO, Newton A. et al. O que explica o desempenho do Brasil no PISA 2015? \textit{Revista Brasileira de Economia}, São Paulo, v. 74, n. 2, p. 167--196, abr. 2020. DOI: 10.5935/0034-7140.20200010. Disponível em: \url{https://periodicos.fgv.br/rbe/article/view/78011}.

\bibitem{OECD2019}
ORGANIZAÇÃO PARA A COOPERAÇÃO E DESENVOLVIMENTO ECONÔMICOS. \textit{PISA 2018 Results (Volume I): What Students Know and Can Do}. Paris: OECD Publishing, 2019. DOI: 10.1787/5f07c754-en. Disponível em: \url{https://www.oecd.org/pisa/}.

\bibitem{RIBEIRO2025}
RIBEIRO JÚNIOR, Marcos dos Santos; LEITE, Bruno Santos. A presença das metodologias ativas nas áreas do conhecimento do CNPq: uma revisão das produções das áreas. \textit{Revista Eletrônica Pesquiseduca}, [S. l.], v. 17, n. 44, p. 97--113, 2025. DOI: 10.58422/repesq.2025.e1751.

% ----------------------------------------------------------------
% Desigualdade e Inclusão Educacional
% ----------------------------------------------------------------

\bibitem{CALLEGARIO2025}
CALLEGARIO ZACCHI, Raquel; TEIXEIRA DE MENEZES, Daniel; GOMES NEY, Marlon. Desigualdades sociais e desempenho escolar no Brasil: o que revelam os microdados do ENEM 2021? \textit{Latitude}, Maceió, v. 18, n. 2, p. 33--59, 2025. DOI: 10.28998/lte.2024.n.2.17872. Disponível em: \url{https://doi.org/10.28998/lte.2024.n.2.17872}.

\bibitem{CAPRARA2020}
CAPRARA, Bernardo Mattes. Condição de classe e desempenho educacional no Brasil. \textit{Educação e Realidade}, Porto Alegre, v. 45, n. 4, 2020. DOI: 10.1590/2175-623693008.

\bibitem{MONTALVAO2014}
MONT'ALVÃO NETO, Arnaldo. Tendências das desigualdades de acesso ao ensino superior no Brasil: 1982--2010. \textit{Educação \& Sociedade}, Campinas, v. 35, n. 127, p. 417--441, 2014. DOI: 10.1590/S0101-73302014000200005.

\bibitem{ZACCHI2016}
ZACCHI, Raquel Callegario; NEY, Marlon Gomes; PONCIANO, Niraldo José. Desigualdades educacionais na educação básica: uma investigação a partir do Exame Nacional do Ensino Médio. \textit{Revista Vértices}, Campos dos Goytacazes, v. 18, n. 1, p. 79--108, 2016. DOI: 10.19180/1809-2667.v18n116-05.

% ----------------------------------------------------------------
% Formação Docente
% ----------------------------------------------------------------

\bibitem{CANTANHEDE2026}
CANTANHEDE, Célia; NAVARRO, Aniela Lopes; MORTALE, Luís Alexandre. Gestão pedagógica baseada no ciclo PDCA na formação docente em metodologias ativas e seus impactos na aprendizagem. \textit{RECIMA21}, [S. l.], v. 7, n. 4, p. e747525, 2026. DOI: 10.47820/recima21.v7i4.7525.

\bibitem{LIBANEO2006}
LIBÂNEO, José Carlos. Reflexividade e formação de professores: outra oscilação do pensamento pedagógico brasileiro. In: PIMENTA, Selma Garrido; GHEDIN, Evandro (org.). \textit{Professor reflexivo no Brasil: gênese e crítica de um conceito}. São Paulo: Cortez, 2006. Disponível em: \url{https://www.worldcat.org/title/didatica/oclc/48964035}.

\bibitem{MIRANDA2022}
MIRANDA, Adriana Tereza da Silva et al. Importância do uso das metodologias ativas para a formação docente. \textit{Brazilian Journal of Development}, [S. l.], v. 8, n. 4, p. 28169--28182, 2022. DOI: 10.34117/bjdv8n4-353.

\bibitem{OLIVEIRA2023}
OLIVEIRA, Francisco Lindoval; NÓBREGA, Luciano; CAVALCANTE, Marcele Alves dos Santos. O uso das metodologias ativas de aprendizagem na formação do professor: das universidades para a prática nas escolas. \textit{Revista Educação Pública}, Rio de Janeiro, v. 23, n. 8, 7 mar. 2023. Disponível em: \url{https://educacaopublica.cecierj.edu.br/artigos/23/8/o-uso-das-metodologias-ativas-de-aprendizagem-na-formacao-do-professor-das-universidades-para-a-pratica-nas-escolas}.

\bibitem{RIOS2026}
RIOS, Pedro Paulo Souza. Didática e metodologias ativas na formação docente: reflexões acerca do fazer pedagógico. \textit{Práxis Educacional}, Vitória da Conquista, v. 22, n. 53, p. e17193, 2026. DOI: 10.22481/praxisedu.v22i53.17193. Disponível em: \url{https://periodicos2.uesb.br/index.php/praxis/article/view/17193}.

\bibitem{SOARES2021}
SOARES, Renato Gomes; CORRÊA, Sandro Luís Pereira; FOLMER, Valter; COPETTI, Jaqueline. A problematização como ferramenta de formação de professores sobre metodologias ativas. \textit{Acta Scientiarum. Education}, Maringá, v. 44, n. 1, p. e52168, 2021. DOI: 10.4025/actascieduc.v44i1.52168. Disponível em: \url{https://periodicos.uem.br/ojs/index.php/ActaSciEduc/article/view/52168}.

\bibitem{VEIGA2014}
VEIGA, Ilma Passos Alencastro. Formação de professores para a educação superior e a diversidade da docência. \textit{Revista Diálogo Educacional}, Curitiba, v. 14, n. 42, p. 327--342, 2014. DOI: 10.7213/dialogo.educ.14.042.DS01.

% ----------------------------------------------------------------
% Tecnologia Educacional
% ----------------------------------------------------------------. Disponível em: \url{https://www.worldcat.org/title/educacao-basica-no-brasil-conceitos-e-praticas/oclc/883601054}.

\bibitem{MEZZARI2011}
MEZZARI, Andréa. O uso da aprendizagem baseada em problemas (ABP) como reforço ao ensino presencial utilizando o ambiente de aprendizagem Moodle. \textit{Revista Brasileira de Educação Médica}, Rio de Janeiro, v. 35, n. 1, p. 114--121, jan./mar. 2011. DOI: 10.1590/S0100-55022011000100015. Disponível em: \url{https://www.scielo.br/j/rbem/a/BLMXqL3Zp8JLmhdRRCn6QCJ/?lang=pt}.

\bibitem{SELWYN2022}
SELWYN, Neil. \textit{Education and technology: Key issues and debates}. 3. ed. London: Bloomsbury Academic, 2022. ISBN: 978-1-3501-4555-9.

% ----------------------------------------------------------------
% Aprendizagem Significativa
% ----------------------------------------------------------------. Disponível em: \url{https://www.bloomsbury.com/uk/education-and-technology-9781350139183/}.

\bibitem{AUSUBEL2000}
AUSUBEL, David P. \textit{The acquisition and retention of knowledge: A cognitive view}. Dordrecht: Kluwer Academic Publishers, 2000. DOI: 10.1007/978-94-015-9454-7. Disponível em: \url{https://www.worldcat.org/title/aquisicao-e-retencao-de-conhecimentos-uma-perspectiva-cognitiva/oclc/48731696}.

\bibitem{NOVAK2008}
NOVAK, Joseph D.; CAÑAS, Alberto J. The theory underlying concept maps and how to construct and use them. Technical Report IHMC CmapTools 2006-01 Rev 01-2008, Florida Institute for Human and Machine Cognition, 2008. Disponível em: \url{http://cmap.ihmc.us/Publications/ResearchPapers/TheoryUnderlyingConceptMaps.pdf}.

% ----------------------------------------------------------------
% Avaliação da Aprendizagem
% ----------------------------------------------------------------

\bibitem{BLACK1998}
BLACK, Paul; WILIAM, Dylan. Assessment and classroom learning. \textit{Assessment in Education: Principles, Policy \& Practice}, Abingdon, v. 5, n. 1, p. 7--74, 1998. DOI: 10.1080/0969595980050102.

\bibitem{SOARES2013}
SOARES, José Francisco; XAVIER, Flávia Pereira. Pressupostos educacionais e estatísticos do Ideb. \textit{Educação \& Sociedade}, Campinas, v. 34, n. 124, p. 903--923, 2013. DOI: 10.1590/S0101-73302013000300013. Disponível em: \url{https://doi.org/10.1590/S0101-73302013000300013}.





% ----------------------------------------------------------------
% Metodologia da Pesquisa
% ----------------------------------------------------------------

\bibitem{BARDIN2016}
BARDIN, Laurence. \textit{Análise de conteúdo}. São Paulo: Edições 70, 2016. Disponível em: \url{https://www.worldcat.org/title/analise-de-conteudo/oclc/973624127}.

\bibitem{CRESWELL2018}
CRESWELL, John William; CRESWELL, J. David. \textit{Research design: Qualitative, quantitative, and mixed methods approaches}. 5. ed. Thousand Oaks, CA: SAGE Publications, 2018. Disponível em: \url{https://www.worldcat.org/title/projeto-de-pesquisa-metodos-qualitativo-quantitativo-e-misto/oclc/1085179213}.

% ----------------------------------------------------------------
% Porter e Cadeia de Valor
% ----------------------------------------------------------------

\bibitem{PORTER1980}
PORTER, Michael E. \textit{Competitive strategy: Techniques for analyzing industries and competitors}. New York: Free Press, 1980.

% ----------------------------------------------------------------
% Normas Técnicas
% ----------------------------------------------------------------. Disponível em: \url{https://www.worldcat.org/title/competitive-strategy-techniques-for-analyzing-industries-and-competitors/oclc/610796358}.

\bibitem{ABNT2018}
ASSOCIAÇÃO BRASILEIRA DE NORMAS TÉCNICAS. ABNT NBR 6023:2018 -- Documentos acadêmicos: referências. Rio de Janeiro: ABNT, 2018.

% ----------------------------------------------------------------
% Produções Adicionais Sobre o Tema
% ----------------------------------------------------------------

\bibitem{ALMEIDA2024}
ALMEIDA, Wellington Augusto Oliveira de et al. A utilização de tecnologias educacionais na aprendizagem baseada em problemas: uma revisão integrativa com ênfase no ensino em saúde. \textit{Caderno Pedagógico}, [S. l.], v. 21, n. 9, p. e8096, 2024. DOI: 10.54033/cadpedv21n9-232.

\bibitem{FREITAS2024}
FREITAS, Vanessa da Silva et al. Formação docente e as metodologias ativas no ensino superior. \textit{REI -- Revista de Educação do UNIDEAU}, [S. l.], v. 4, n. 1, p. e170, 2024. Disponível em: \url{https://periodicos.ideau.com.br/index.php/rei/article/view/170}.

\bibitem{MACHADO2025}
MACHADO, Ariane et al. Formação docente em metodologias ativas: lições aprendidas. \textit{PRATICA -- Revista Multimédia de Investigação em Inovação Pedagógica e Práticas de e-Learning}, [S. l.], v. 8, n. 3, p. 1--11, 2025. DOI: 10.34630/pel.v8i3.6360.

\bibitem{OLIVEIRA2023B}
OLIVEIRA, L. L. de; MOURA, N. P. R.; TANAJURA, D. M. Aprendizagem baseada em problemas e o currículo tradicional na educação em enfermagem: uma revisão bibliográfica. \textit{Educationis}, [S. l.], v. 3, n. 1, p. 34--41, 2015. DOI: 10.6008/SPC2318-3047.2015.001.0004.

% ============================================================
% Fim do Capítulo 6 -- Referências Bibliográficas
% ============================================================

\end{thebibliography}
